% !TEX TS-program = lualatex
% !TEX encoding = UTF-8 Unicode
% What's New のページのスクリーンショットのための見本。
\documentclass[11pt,a4paper]{ltjsarticle}
\usepackage[haranoaji]{luatexja-preset}
\usepackage{amsmath,amssymb,amsthm}
\usepackage{pgfplots}
\pgfplotsset{compat=1.18}
\usepackage[margin=23mm]{geometry}
\usepackage{microtype}
\usepackage[hidelinks]{hyperref}
\hypersetup{pdftitle={波からフーリエ級数へ}}
\setlength{\parskip}{0.45em}
\newtheorem{theorem}{定理}
\theoremstyle{definition}
\newtheorem{definition}{定義}

\begin{document}
\begin{center}
  {\LARGE 波からフーリエ級数へ}\\[0.5em]
  {\small フーリエ級数入門}
\end{center}

\section{波を重ねて関数を表す}\label{sec:waves}
複雑な信号も、単純な振動の重ね合わせとして捉えることができます。
周期 $2\pi$ の関数に対するフーリエ級数は、次の形で表されます。
\begin{equation}\label{eq:series}
  f(x) \sim \frac{a_0}{2}
  + \sum_{n=1}^{\infty}
    \bigl(a_n\cos nx + b_n\sin nx\bigr),
\end{equation}
ここで、各成分の係数は次の積分で求められます。
\begin{equation}\label{eq:coefficients}
  a_n = \frac{1}{\pi}\int_{-\pi}^{\pi}f(x)\cos nx\,dx,
  \qquad
  b_n = \frac{1}{\pi}\int_{-\pi}^{\pi}f(x)\sin nx\,dx.
\end{equation}
級数がどこで何に収束するかは、第\ref{sec:convergence}節の
定理\ref{thm:dirichlet}が答えます。

\begin{definition}\label{def:piecewise}
区間 $[-\pi,\pi]$ を有限個の小区間に分け、各小区間で $f$ と $f'$ が
連続で、端で片側極限をもつとき、$f$ は\emph{区分的に滑らか}であるといいます。
\end{definition}

\section{矩形波を近似する}\label{sec:square}
区間 $(-\pi,0)$ で $f(x)=-1$、区間 $(0,\pi)$ で $f(x)=1$ とする
奇関数の矩形波では、余弦の項がすべて消え、
\eqref{eq:coefficients} から奇数の $n$ について $b_n=4/(n\pi)$ となります。
\eqref{eq:series} の最初の $N$ 個の非零項を用いると、次の近似が得られます。
\begin{equation}\label{eq:partial}
  S_N(x)=\frac{4}{\pi}\sum_{k=0}^{N-1}
  \frac{\sin\bigl((2k+1)x\bigr)}{2k+1}.
\end{equation}

\begin{figure}[ht]
\centering
\begin{tikzpicture}
\begin{axis}[
  width=0.93\linewidth, height=6.3cm,
  axis lines=middle, xlabel={$x$}, ylabel={$y$},
  xmin=-3.5, xmax=3.5, ymin=-1.5, ymax=1.5,
  xtick={-3.14159265,0,3.14159265},
  xticklabels={$-\pi$,$0$,$\pi$}, ytick={-1,1},
  samples=401, domain=-pi:pi,
  legend style={at={(0.5,-0.15)}, anchor=north,
    draw=none, legend columns=2},
  every axis plot/.append style={thick}
]
\addplot[blue!75!black] {4/pi*sin(deg(x))};
\addlegendentry{1項の近似}
\addplot[red!80!black]
  {4/pi*(sin(deg(x))+sin(deg(3*x))/3+sin(deg(5*x))/5
        +sin(deg(7*x))/7+sin(deg(9*x))/9)};
\addlegendentry{5項の近似}
\end{axis}
\end{tikzpicture}
\caption{矩形波の部分和 \eqref{eq:partial}}\label{fig:square}
\end{figure}

図\ref{fig:square}のように、項数を増やすと、
不連続点から離れた場所では近似が改善します。

\clearpage
\section{収束}\label{sec:convergence}
級数 \eqref{eq:series} がどこで何に収束するかは、次の古典的な結果が教えてくれます。

\begin{theorem}[ディリクレ]\label{thm:dirichlet}
$f$ は周期 $2\pi$ で、定義\ref{def:piecewise}の意味で区分的に滑らかとします。
このとき、フーリエ級数 \eqref{eq:series} はすべての点 $x$ で収束し、その和は
\[
  \frac{f(x^{+})+f(x^{-})}{2}
\]
です。特に、$f$ が連続な点では、級数は $f(x)$ に収束します。
\end{theorem}

第\ref{sec:square}節の矩形波では、$x=0$ での和は $-1$ と $1$ の平均の $0$ です。

\section{ギブズ現象}\label{sec:gibbs}
不連続点の近くでは、項数を増やしても一定の行き過ぎが残り、
$x=0$ の近くの $S_N$ の最大値は、$N\to\infty$ のとき
\begin{equation}\label{eq:gibbs}
  \frac{2}{\pi}\int_0^{\pi}\frac{\sin t}{t}\,dt \approx 1.1790
\end{equation}
に近づきます。これが\textbf{ギブズ現象}です。これは定理\ref{thm:dirichlet}と
矛盾しません。$x\neq0$ を固定すれば、部分和 \eqref{eq:partial} は
$f(x)$ に収束しますが、行き過ぎの場所が不連続点に近づいていくのです。

\begin{theorem}\label{thm:mean}
$f$ が2乗可積分ならば、\eqref{eq:series} の部分和は $f$ に平均収束します。
\[
  \lim_{N\to\infty}\int_{-\pi}^{\pi}\bigl|f(x)-S_N(x)\bigr|^2\,dx=0.
\]
\end{theorem}

定理\ref{thm:mean}と値 \eqref{eq:gibbs} から、図\ref{fig:square}の行き過ぎは、
高さは減らないものの、占める区間はしだいに狭くなることがわかります。
なお、本文に全角で「Ｓ＿Ｎ（ｘ）＝１」と書くと、数式にはなりません。
\end{document}
